% =============================================================================
% Kapitel 2: Notwendigkeit und Anwendungsfelder von Off-Grid Schutzsensorik
% =============================================================================

\chapter{Notwendigkeit und Anwendungsfelder von Off-Grid Schutzsensorik}
\label{chap:motivation}

% TODO: update intro when done proofreading
Dieses Kapitel legt die anwendungsseitige Grundlage der Arbeit. Es wird zunächst
dargelegt, warum konventionelle Funkkommunikationstechnologien in Gefahrgebieten
strukturell versagen und welche Anforderungen sich daraus an ein autonomes
Kommunikationssystem ergeben. Anschließend werden vier repräsentative
Anwendungsfelder analysiert, in denen diese Anforderungen auftreten. Das Kapitel
schließt mit einer tabellarischen Gegenüberstellung der Anwendungsfelder, aus der
die technischen Designziele für das in Kapitel~\ref{chap:hardware} entwickelte
System abgeleitet werden.


% -----------------------------------------------------------------------------
\section{Versagen konventioneller Kommunikationsinfrastruktur in Gefahrgebieten}
\label{sec:komm-versagen}
% -----------------------------------------------------------------------------

Moderne drahtlose Kommunikationstechnologien wie WLAN, Bluetooth und % TODO: def acro
Mobilfunknetze der vierten und fünften Generation sind für den Einsatz in
erschlossenen, infrastrukturell versorgten Umgebungen konzipiert. In
Gefahrgebieten -- definiert als Orte, an denen Personen physischen, chemischen
oder biologischen Risiken ausgesetzt sind und an denen gleichzeitig keine oder
nur unzuverlässige Kommunikationsinfrastruktur vorhanden ist -- stoßen diese
Technologien an fundamentale Grenzen.

WLAN nach IEEE 802.11 ist auf Reichweiten von typischerweise 30 bis
\SI{100}{\m} im Innenbereich begrenzt und erfordert die Installation von
Access-Points in regelmäßigen Abständen. Dies setzt eine stabile,
netzgespeiste Infrastruktur voraus, die in Katastrophengebieten, im Untertagebergbau
oder in abgelegenen Industrieanlagen häufig nicht vorhanden ist. Bluetooth ist
mit einer Reichweite von bis zu \SI{10}{\m} (Klasse 2) für
weiträumige Gefahrgebiete von vornherein ungeeignet. Mobilfunknetze
bieten zwar eine großflächige Abdeckung, weisen jedoch in Katastrophenszenarien
eine systemimmanente Schwäche auf: Basisstationen werden häufig durch dasselbe
Ereignis zerstört oder überlastet, das den Notfall verursacht -- Erdbeben, Überschwemmungen
oder Brände~\cite{sciullo2020locate}. Untertageumgebungen wie Bergwerksstollen
stellen eine besondere Herausforderung dar, da Funksignale durch Gestein stark
gedämpft werden und herkömmliche Mobilfunk- oder WLAN-Signale bereits nach wenigen
Metern vollständig absorbiert werden.

Das Konzept der \textbf{Off-Grid-Kommunikation} bezeichnet drahtlose
Netzwerke, die vollständig autark -- ohne Anbindung an bestehende
Telekommunikationsinfrastruktur -- operieren. Für Schutzsensorsysteme in
Gefahrgebieten ergeben sich daraus drei grundlegende Anforderungen, die in
Tabelle~\ref{tab:komm_anforderungen} zusammengefasst sind.

\begin{table}[htbp]
    \centering
    \caption{Grundlegende Kommunikationsanforderungen für Off-Grid Schutzsensorik}
    \label{tab:komm_anforderungen}
    \begin{tabularx}{\textwidth}{l X}
        \toprule
        \textbf{Anforderung} & \textbf{Begründung} \\
        \midrule
        Infrastrukturunabhängigkeit &
        Das Netzwerk muss ohne externe Basisstationen, Internetanbindung
        oder Netzspannung operieren. \\
        \addlinespace
        Energieautarkie &
        Endgeräte müssen über akkubetriebene oder energy-harvesting-gestützte
        Quellen dauerhaft betrieben werden können, da Netzstrom im Einsatzgebiet
        nicht verfügbar ist. \\
        \addlinespace
        Reichweite und Gebäudedurchdringung &
        Das System muss Distanzen von mehreren hundert Metern bis zu einigen
        Kilometern überbrücken und dabei Hindernisse wie Betonwände,
        Tunnelwände oder Waldvegetation durchdringen. \\
        \bottomrule
    \end{tabularx}
\end{table}

LoRaWAN auf Basis der LoRa"=Modulationstechnologie adressiert alle drei
Anforderungen: Es operiert im lizenzfreien ISM"=Band, erreicht Reichweiten
von 2 bis \SI{5}{\km} in städtischen und über \SI{10}{\km} in ländlichen
Umgebungen, benötigt im Sendebetrieb lediglich wenige Milliampere und kann
vollständig ohne Internetanbindung in einem privaten, autarken Netzwerk betrieben
werden. Die folgenden Abschnitte zeigen anhand konkreter Anwendungsfelder, in welchen Szenarien diese 
Eigenschaften entscheidend sind.

% -----------------------------------------------------------------------------
\section{Use Case 1: Industrie und Bergbau}
\label{sec:uc-industrie}
% -----------------------------------------------------------------------------

\subsection{Gefährdungslage}

Der Untertagbergbau gehört zu den gefährlichsten Arbeitsbereichen weltweit.
\citeauthor{bls_mining2023} berichten, dass die Todesrate in der US-amerikanischen
Kohleindustrie mit \num{19,6} Todesfällen pro \num{100000} Vollzeitbeschäftigten
im Jahr 2021 mehr als das Fünffache des Durchschnittswerts der Privatwirtschaft
(\num{3,8}) beträgt~\cite{bls_mining2023}. Als häufigste Unfallursachen identifiziert
eine Analyse chinesischer Bergwerksunfälle zwischen 2017 und 2022 Dacheinbrüche
sowie Gasexplosionen~\cite{sciencedirect_coal2024}. Methan und
Kohlenmonoxid akkumulieren in schlecht belüfteten Stollen und können
bei Erreichen der unteren Explosionsgrenze -- für Methan bei \SI{5}{\percent}
Volumenanteil in Luft -- zu verheerenden Explosionen führen~\cite{wiki_mining}.

In chemischen Anlagen und Raffinerien stellt die ATEX"=Richtlinie 2014/34/EU
den regulatorischen Rahmen für Geräte in explosionsgefährdeten Bereichen dar.
Diese Umgebungen sind in Zonen unterteilt (Zone 0 bis 2 für Gase, Zone 20 bis 22
für Stäube), die Anforderungen an die Zündschutzart von elektrischen Betriebsmitteln
stellen. Ein verteiltes Gasmesssystem muss daher nicht nur funktional zuverlässig,
sondern auch normkonform ausgelegt sein.

\subsection{Kommunikationsherausforderung}

Im Untertagebergbau fehlt typischerweise jede Form von Mobilfunkabdeckung.
Funksignale werden durch Gestein stark gedämpft; die Signaldämpfung in Stollen
kann je nach Gesteinsart und Geometrie mehrere zehn Dezibel pro \SI{100}{\m}
betragen~\cite{niosh_wsn}. Bestehende kabelgebundene Systeme sind in ihrer
Skalierbarkeit eingeschränkt und bieten keine Mobilität der Endgeräte. Gleichzeitig
ist eine zeitkritische Alarmübertragung erforderlich: Im Falle einer
Gaskonzentration oberhalb des Schwellwerts muss ein Warnsignal innerhalb von
Sekunden die gesamte betroffene Belegschaft erreichen.

\subsection{LoRaWAN als Lösungsansatz}

\citeauthor{baert2018lorawan} demonstrieren in ihrer Arbeit zu industriellen
IoT"=Anwendungen mit LoRaWAN, dass LoRa"=Signale aufgrund der CSS"=Modulation
auch in Umgebungen mit starker Mehrwegeausbreitung und
Signaldämpfung -- wie sie in Tunneln und Stollen auftreten -- noch zuverlässig
demoduliert werden können~\cite{baert2018lorawan}. \citeauthor{smartmininghelmet}
berichten in einer Prototypstudie, dass LoRa"=Kommunikation in simulierten
Untertagebedingungen eine zuverlässige Datenübertragung bis zu \SI{300}{\m}
bei einem Stromverbrauch ermöglichte, der einen ununterbrochenen Betrieb von
über 12 Stunden erlaubt~\cite{smartmininghelmet}.

Ein LoRaWAN"=basiertes Schutzsensorsystem für dieses Anwendungsfeld würde aus
mobilen Endgeräten (Gasdetektoren, am Körper oder Helm getragen), einem oder
mehreren portablen Gateways sowie einem lokal betriebenen Network-Server bestehen.
Die mobilen Endgeräte erfassen kontinuierlich Gaskonzentrationen und übertragen
bei Überschreitung eines Schwellwerts ein Alarmpaket. Da ADR für mobile % TODO: define ADR acronym
Geräte nicht geeignet ist (vgl. Abschnitt~\ref{subsec:adr}), muss der Spreading
Factor entweder manuell konfiguriert oder durch ein adaptives, mobilitätsbewusstes
Verfahren gesteuert werden.

Hinsichtlich der ATEX"=Konformität ist zu beachten, dass der in dieser Arbeit
entwickelte Prototyp auf Basis des STM32WL55 nicht für den Einsatz in
explosionsgefährdeten Bereichen zertifiziert ist. Die ATEX"=Zertifizierung stellt
einen separaten Schritt dar, der für eine Serienproduktion durch
Descap GmbH erforderlich wäre.

% -----------------------------------------------------------------------------
\section{Use Case 2: Katastrophenschutz und Feuerwehr}
\label{sec:uc-katastrophenschutz}
% -----------------------------------------------------------------------------

\subsection{Gefährdungslage}

Einsatzkräfte des Katastrophenschutzes und der Feuerwehr operieren per Definition
in Umgebungen, in denen die Infrastruktur entweder nicht vorhanden oder durch
das Ereignis selbst zerstört worden ist. Laut NFPA kamen im Jahr 2023
in den USA 89 Feuerwehrleute im Dienst ums Leben -- der zweithöchste Wert seit
2013~\cite{nfpa2023}. Als besonders gefährlich gelten Waldbrände: Einsatzkräfte
können sich in sich schnell verändernden Brandgebieten ohne Sichtverbindung und
ohne Mobilfunknetz befinden. Der Verlust der Kommunikation zu Teamkollegen
oder zur Einsatzleitung ist in diesen Szenarien lebensbedrohlich.

Bei Einsturzszenarien nach Erdbeben oder Gebäudeeinstürzen müssen Such- und
Rettungsteams in eingestürzten Gebäudestrukturen operieren, in denen
Mobilfunksignale durch Stahlbeton vollständig abgeschirmt werden. Auch
Überflutungsgebiete erfordern eine vom terrestrischen Netz unabhängige
Kommunikation, da Basisstationen und Stromversorgung häufig vom selben Ereignis
betroffen sind wie die Einsatzkräfte.

\subsection{Kommunikationsherausforderung}

Die Anforderungen an Kommunikationssysteme in Katastrophenschutzszenarien
unterscheiden sich in einem wesentlichen Punkt von industriellen Anwendungen:
Das System muss \textbf{rasch deployierbar} sein. Im Gegensatz zu fest
installierten Industriesystemen muss ein mobiles LoRaWAN"=Netz innerhalb von
Minuten nach Eintreffen am Einsatzort betriebsbereit sein. Dies setzt voraus,
dass das Gateway batteriebetrieben, kompakt und von nicht"=technischem Personal
in Betrieb zu nehmen ist.

Darüber hinaus ist im Katastrophenschutz die Personenverfolgung (\textit{Personnel
Tracking}) eine zentrale Funktion: Die Einsatzleitung muss jederzeit wissen,
wo sich einzelne Trupps befinden. LoRa unterstützt nativ eine
TDOA"=basierte (\textit{Time Difference of Arrival}) Ortung, die ohne
GPS"=Empfänger am Endgerät auskommt und damit für Innenraumeinsätze
relevant ist~\cite{semtech_lora_basics}.

\subsection{LoRaWAN als Lösungsansatz}

\citeauthor{sciullo2020locate} präsentieren mit \textit{LOCATE} ein
LoRaWAN"=basiertes mobiles Einsatzmanagementsystem, das explizit für
Katastrophenschutzszenarien ohne bestehende Infrastruktur entwickelt wurde.
Das System nutzt eine Mesh"=Topologie, bei der einzelne Einsatzkräfte als
Relaisknoten fungieren und damit die Netzabdeckung dynamisch in das
Einsatzgebiet hinein erweitern~\cite{sciullo2020locate}. Dies ist ein
wesentlicher Vorteil gegenüber der klassischen Stern"=Topologie, bei der alle
Endgeräte in direkter Reichweite eines Gateways sein müssen.

Für Waldbrandeinsätze bietet die Doppler"=Resistenz von LoRa (vgl.
Abschnitt~\ref{sec:lora-physical}) einen praktischen Vorteil: Auch bei
bewegten Einsatzkräften bleibt die Verbindungsqualität stabil, ohne dass
eine Frequenznachführung erforderlich ist.


% -----------------------------------------------------------------------------
\section{Use Case 3: Militär und Sicherheit}
\label{sec:uc-militar}
% -----------------------------------------------------------------------------

\subsection{Gefährdungslage}

Militärische und sicherheitsbehördliche Operationen in abgelegenen oder
feindlich kontrollierten Gebieten stellen besondere Anforderungen an
Kommunikationssysteme. EOD"=Teams, die mit der Entschärfung von % TODO: def acro
Sprengvorrichtungen beauftragt sind, operieren häufig in Umgebungen, die
gezielt elektromagnetisch gestört werden (\textit{Electronic Warfare},
Jamming). Grenzüberwachungsmissionen in infrastrukturschwachen Regionen --
etwa in Wüsten- oder Gebirgsterrain -- können sich über hunderte von Kilometern
erstrecken, ohne dass eine Mobilfunkabdeckung vorhanden ist.
UN"=Friedensmissionen in Entwicklungsländern sehen sich mit dem
vollständigen Fehlen ziviler Telekommunikationsinfrastruktur konfrontiert.

\subsection{Kommunikationsherausforderung}

Klassische Militärfunksysteme wie VHF"=Handfunkgeräte sind zuverlässig, % TODO: def acro
aber schwer, energiehungrig und durch ihre charakteristischen Emissionen
detektierbar. Für sicherheitskritische IoT"=Anwendungen -- etwa ein verteiltes
Netzwerk von Bewegungs- oder Schadstoffsensoren entlang einer
Überwachungslinie -- sind dedizierte Funkgeräte weder skalierbar noch praktikabel.

LoRaWAN bietet in diesem Kontext zwei relevante Eigenschaften: Erstens
ist der \textit{RF Footprint} eines LoRa"=Endgeräts aufgrund der geringen
Sendeleistung (typischerweise \SIrange{14}{22}{dBm}) und der auf den ersten
Blick rauschartig erscheinenden CSS"=Modulation schwer zu detektieren.
Zweitens ermöglicht die niedrige Datenrate eine sehr stromsparende,
intermittierende Übertragung, die mit kleinen Primärzellen oder solarbetriebenen
Akkus über Monate betrieben werden kann~\cite{morin2017device}.

\subsection{LoRaWAN als Lösungsansatz}

Ein autonomes Sensornetzwerk für Grenz- oder Perimeterschutz könnte aus
solarbetriebenen LoRa"=Knoten bestehen, die in regelmäßigen Abständen entlang
einer Überwachungslinie platziert werden und nur bei Detektion eines Ereignisses
ein Uplink"=Paket senden. Die geringe Duty"=Cycle"=Auslastung (Abschnitt~\ref{subsec:frequenzbaender})
fügt sich dabei gut in dieses ereignisgesteuerte Übertragungsschema ein.
Ein portables Gateway, betrieben von einem Fahrzeug oder einem vorgeschobenen
Posten, sammelt die Sensordaten und leitet sie an eine taktische
Lagezentrale weiter -- ohne Anbindung an öffentliche Kommunikationsnetze.

Es sei darauf hingewiesen, dass militärische Anwendungen spezifische
Anforderungen an Verschlüsselung und Zertifizierung stellen, die über den
im Rahmen dieser Arbeit betrachteten zivilen AES"=128"=Standard hinausgehen.
Der in dieser Arbeit entwickelte Prototyp ist als Basis für zivile
Schutzsensorik konzipiert; eine Erweiterung in Richtung militärischer
Normen (z.B. NATO STANAG) würde weitergehende Untersuchungen erfordern.


% -----------------------------------------------------------------------------
\section{Use Case 4: Umweltmonitoring in Extremgebieten}
\label{sec:uc-umwelt}
% -----------------------------------------------------------------------------

\subsection{Gefährdungslage}

Umweltmessstationen in Extremlagen -- darunter Vulkane, Gletscher,
Hochwasser"=Frühwarnsysteme in Entwicklungsländern und Waldbrandfrühwarnsysteme --
stehen vor einem grundlegenden Widerspruch: Die Gebiete, in denen Messungen am
dringendsten benötigt werden, sind häufig diejenigen, die am schlechtesten
erschlossen sind. Vulkane erfordern Sensoren in unmittelbarer Nähe der
Ausbruchszone, wo keine Infrastruktur existiert und physischer Zugang
für Menschen lebensgefährlich ist. Gletschermessstationen in arktischen oder
hochalpinen Regionen müssen über Monate hinweg -- teils durch polare Nacht
mit eingeschränkter Solarenergie -- autonom operieren.

Hochwasser"=Frühwarnsysteme in Entwicklungsländern sind auf einfach
installierbare, kostengünstige und wartungsarme Sensorknoten angewiesen, die
keine telekommunikationstechnische Infrastruktur voraussetzen. Das Fehlen
eines solchen Systems hatte in der Vergangenheit wiederholt katastrophale
Folgen: Frühwarnsysteme können die Vorlaufzeit für Evakuierungen von
Minuten auf Stunden verlängern~\cite{adeel2019wsn}.

\subsection{Kommunikationsherausforderung}

Die dominante Herausforderung in diesem Anwendungsfeld ist
\textbf{Energieautarkie über lange Zeiträume}. Ein Sensorknoten, der an einem
Vulkanhang oder in einem Flussbett installiert ist, kann nicht regelmäßig
für einen Batteriewechsel aufgesucht werden. Das System muss daher entweder
mit Primärbatterien über mehrere Jahre oder mit solarbetriebenen Akkus
betrieben werden. Dies erfordert einen durchschnittlichen Stromverbrauch im
Bereich weniger Mikroampere -- ein Wert, der nur durch konsequente Nutzung
von Schlafmodi und sehr niedrige Sendefrequenz erreichbar ist.

Gleichzeitig gibt es in diesen Gebieten keine Mobilfunkabdeckung und keine
Möglichkeit, einen kabelgebundenen Network-Server zu betreiben. Das Gateway
muss daher ebenfalls solar- oder batteriebetrieben und wetterfest sein.

\subsection{LoRaWAN als Lösungsansatz}

\citeauthor{lorawan_aqms2022} zeigen, dass ein LoRaWAN"=basiertes
Luftqualitätsmesssystem mit solarbetriebenen Knoten einen dauerhaften,
netzunabhängigen Betrieb ermöglicht~\cite{lorawan_aqms2022}. Die Verwendung von
SF12 -- dem Spreading Factor mit der höchsten Empfängerempfindlichkeit von
bis zu \SI{-148}{dBm} -- ermöglicht Reichweiten von mehreren Kilometern auch
bei ungünstigen Ausbreitungsbedingungen, wie sie in Berggebieten durch
Abschattung auftreten. Class-A"=Betrieb mit langen Schlafphasen zwischen den
Messzyklen minimiert den Energieverbrauch auf ein Niveau, das mit kleinen
Solarpanelen vereinbar ist.

Ein realistisches Systemszenario für dieses Anwendungsfeld könnte einen Knoten
vorsehen, der alle \SI{5}{\min} eine aggregierte Messung (Minimum, Mittelwert,
Maximum) überträgt~\cite{semtech_lora_basics}. Bei SF12, BW~\SI{125}{kHz}
und einer Payload von 12~Bytes beträgt die Time"=on"=Air pro Übertragung
ca. \SI{1318}{ms}, was bei einer Wiederholrate von \SI{5}{\min} einer
Duty"=Cycle"=Auslastung von \SI{0.44}{\percent} entspricht -- unterhalb der
in Europa zulässigen 1\,\%.


% -----------------------------------------------------------------------------
\section{Vergleich und Ableitung der Systemanforderungen}
\label{sec:anforderungsableitung}
% -----------------------------------------------------------------------------

Die vier analysierten Anwendungsfelder unterscheiden sich in ihren konkreten
Rahmenbedingungen erheblich, teilen jedoch eine gemeinsame Grundstruktur:
fehlende Infrastruktur, strenge Energiebudgets und die Notwendigkeit einer
zuverlässigen Datenübertragung über nicht-triviale Distanzen. Tabelle~\ref{tab:usecases}
fasst die wesentlichen Merkmale der Anwendungsfelder zusammen.

\begin{table}[htbp]
    \centering
    \caption{Vergleich der Anwendungsfelder hinsichtlich Systemanforderungen}
    \label{tab:usecases}
    \begin{tabularx}{\textwidth}{l X X X X}
        \toprule
        & \textbf{Industrie /} & \textbf{Katastrophen-} & \textbf{Militär /}
        & \textbf{Umwelt-} \\
        & \textbf{Bergbau} & \textbf{schutz} & \textbf{Sicherheit}
        & \textbf{monitoring} \\
        \midrule
        Infrastruktur & Keine (Untertagebau) & Zerstört / keine & Keine / feindlich
        & Keine \\
        \addlinespace
        Reichweite & \SIrange{100}{500}{\m} & \SIrange{500}{2000}{\m}
        & \SIrange{1}{10}{\km} & \SIrange{2}{15}{\km} \\
        \addlinespace
        Batterielaufzeit & 8--24\,h & 4--48\,h & Monate & Monate--Jahre \\
        \addlinespace
        Latenztoleranz & Niedrig ($<$\,\SI{5}{\s}) & Mittel ($<$\,\SI{30}{\s})
        & Hoch ($<$\,\SI{5}{\min}) & Sehr hoch ($<$\,\SI{30}{\min}) \\
        \addlinespace
        Mobilität & Hoch & Hoch & Mittel & Keine \\
        \addlinespace
        Topologie & Stern / Mesh & Mesh & Stern & Stern \\
        \addlinespace
        Normierung & ATEX, DGUV & --- & STANAG (optional) & --- \\
        \bottomrule
    \end{tabularx}
\end{table}

Aus dem Vergleich lassen sich folgende Designziele für das in dieser Arbeit
entwickelte System ableiten:

\begin{itemize}
    \item \textbf{Autonomer Betrieb:} Das System muss ohne Internetanbindung
    und ohne externe Netzspannung vollständig funktionsfähig sein. Als
    Network"=Server kommt ChirpStack auf einem Einplatinenrechner (Raspberry Pi)
    zum Einsatz.

    \item \textbf{Energieeffizienz:} Die Sensoreinheit muss im Akkubetrieb
    eine Mindestlaufzeit von 8 Stunden erreichen, um einen Arbeitsschicht-konformen
    Betrieb zu ermöglichen. Dies erfordert die konsequente Nutzung von
    Schlafmodi auf dem STM32WLE5 (vgl. Abschnitt~\ref{sec:energieeffizienz}).

    \item \textbf{Konfigurierbare Reichweite:} Der Spreading Factor muss
    softwareseitig konfigurierbar sein, um den Trade-off zwischen Reichweite
    und Energieverbrauch für das jeweilige Einsatzszenario optimieren zu können.

    \item \textbf{Topologievergleich:} Da Stern- und Mesh"=Topologie je nach
    Anwendungsfeld unterschiedliche Stärken aufweisen, werden beide
    Konfigurationen im experimentellen Teil der Arbeit
    (Kapitel~\ref{chap:evaluation}) untersucht und verglichen.

    \item \textbf{Skalierbarkeit auf Industrieanforderungen:} Der Prototyp soll
    so ausgelegt sein, dass eine spätere Integration in zertifizierungsfähige
    Hardware (ATEX"=Zonen) architekturell möglich ist, auch wenn die
    Zertifizierung selbst nicht Gegenstand dieser Arbeit ist.
\end{itemize}
